QPS is computed as $1/\text{mean latency}$, not measured throughput
under load. This is the deployment question (``serial latency under
no contention'') rather than a saturation curve --- multi-threaded
throughput with batching has its own axis (figure 07, parallel
scaling). The chart conveys exactly the same data as figure 01 with
the x-axis flipped (higher = better); reading both figures
side-by-side is redundant. Use this one when comparing against
QPS-based service-level objectives.

Each point carries asymmetric x-axis error bars: the short positive
arm is the across-rep $95\%$ CI on the mean (precision of the
measurement), and the long negative arm extends to the
$\text{p99-throughput}$ point ($1/\text{p99-latency}$). A wide
negative whisker says ``the mean is achievable on average but the
tail reaches further left'' --- useful for distinguishing schemes
whose per-query latency is consistent (short whisker) from ones
whose tail dominates the worst-case experience (long whisker). Both
arms use the same raw \texttt{latency-us} data; no new sweeps were
needed.
